% Inactive modular snapshot; edit the root neurips_2026.tex instead.
\documentclass{article}
\PassOptionsToPackage{table}{xcolor}
\usepackage{iclr2027_conference,times}
\usepackage[T1]{fontenc}
\usepackage[utf8]{inputenc}
\usepackage{amsmath,amssymb,amsthm,booktabs,multirow,adjustbox}
\usepackage{xcolor}
\usepackage{enumitem,graphicx,microtype,url,hyperref}
\hypersetup{hidelinks}
\newtheorem{proposition}{Proposition}
\newtheorem{corollary}[proposition]{Corollary}
\newcommand{\pre}{\mathrm{pre}}
\newcommand{\off}{\mathrm{off}}
\newcommand{\cross}{\mathrm{cross}}
\newcommand{\diag}{\operatorname{diag}}
\newcommand{\op}[1]{\left\|#1\right\|_2}
\newcommand{\fro}[1]{\left\|#1\right\|_F}
% Primitive input avoids LaTeX file hooks between a tabular row and booktabs rules.
\makeatletter
\newcommand{\tablerows}[1]{\@@input #1 }
\makeatother
\newcommand{\allocationbar}[4]{\makebox[1em][l]{#1}\begingroup\setlength{\unitlength}{\dimexpr\linewidth-1em\relax}\textcolor{black!80}{\rule{#2\unitlength}{1.4ex}}\textcolor{black!45}{\rule{#3\unitlength}{1.4ex}}\textcolor{black!15}{\rule{#4\unitlength}{1.4ex}}\endgroup}
\title{Confinement, Flexibility, and Mixing\\in Spectral Fine-Tuning}
\author{Anonymous authors}
% Keep anonymous; do not enable \iclrfinalcopy for submission.
\begin{document}
\maketitle

\begin{abstract}
Where does parameter-efficient fine-tuning modify a pretrained model? Rank and trainable parameter count do not specify an update's location relative to the pretrained weight structure. We study this location by decomposing updates into leading, tail, and cross-subspace blocks in pretrained singular coordinates. Fixed coefficients, within-tail rotations, and ambient diagonal scalers isolate different forms of adaptation flexibility. Exact approximation identities separate restrictions within the tail from restrictions imposed by its span. Mixing bounds distinguish strict tail confinement from cross-subspace decoupling, including a practical construction with a leading core. In a five-task RoBERTa-base intervention, two-sided mixing regularization reduces cross energy from approximately 43\% to 10--27\%, while the leading block retains 44--64\% of update energy. Leading-subspace drift falls in all nine paired task/seed comparisons, alongside substantial update shrinkage and task-mean score decreases of $0.2$--$1.0$ points. Broader language-understanding and generation evaluations establish the practical viability of these instruments. The study distinguishes three aspects of spectral adaptation: representational freedom, interaction with pretrained subspaces, and effective update magnitude.
\end{abstract}
% Inactive modular snapshot; edit the root neurips_2026.tex instead.
\section{Introduction}
\label{sec:intro}

Parameter-efficient fine-tuning (PEFT) adapts a pretrained model through a restricted set of trainable parameters. Rank, parameter count, and task performance are natural descriptions of these restrictions, but they leave an important question open: where does the effective update lie relative to the structure of the pretrained weights? Two updates can have the same rank and magnitude while modifying different pretrained directions. Likewise, similar downstream performance need not imply similar spectral changes, as studies comparing LoRA and full fine-tuning demonstrate~\citep{hu2022lora,shuttleworth2025illusion}.

The singular value decomposition of a pretrained weight matrix provides a fixed reference for this question. Its leading and tail singular subspaces define a partition of the update into leading--leading, tail--tail, and two cross blocks. This is a geometric distinction, not a semantic assignment: a large singular value does not by itself identify stored knowledge, and a small singular value does not identify dispensable noise. The reference frame instead lets us ask what a constraint preserves, what freedom it removes, and how a relaxation interacts with the original directions.

This reference frame is already used by spectral fine-tuning methods, including singular-value tuning, singular-vector coefficient tuning, and regularized adaptation of singular factors~\citep{han2023svdiff,lingam2024svft,cao2025sorsa}. Our question is not whether pretrained spectral information can be used for adaptation, but what different uses of that information actually constrain. In particular, orthonormality of trainable factors, confinement to a pretrained subspace, and stability of the largest singular directions are different properties.

We investigate these distinctions through a hierarchy of spectral constraints. Fixed tail coefficients restrict both the available span and the directions expressible inside it. Within-tail rotations relax the second restriction without relaxing the first. Ambient diagonal scaling acts outside the singular coordinate system and can introduce leading--tail interaction. The implementation names LinLoRA, OrthoLoRA, UILinLoRA, and UIOrthoLoRA identify these constructions. We use them to isolate representational freedoms and study regularization of their interactions with the pretrained spectrum.

The analysis yields two separations. First, exact projection identities distinguish an outside-tail approximation residual from the additional error imposed by a diagonal tail core. They hold for any target perturbation, without assuming that task-relevant directions align with the tail. Second, the meaning of mixing control depends on the core being scaled. A tail-only core recovers confinement when mixing vanishes. A practical leading-plus-tail core instead permits direct leading adaptation even when its cross blocks vanish. Thus suppressing mixing need not preserve the leading weight block, and preserving that block need not preserve its rank ordering after adaptation.

We then examine a mixing-regularization intervention across five GLUE tasks. Penalizing both sides reduces cross energy from approximately 43\% to 10--27\% and increases tail allocation. Yet the leading block's share rises in four of five task means. This distinction would be hidden by reporting only aggregate off-tail leakage. Leading-subspace drift also falls in every paired task/seed comparison, accompanied by substantial update shrinkage and task-mean score decreases of $0.2$--$1.0$ points. The experiment establishes a geometry--magnitude--performance trade-off; it does not isolate an effect beyond magnitude-only regularization.

Our contributions are a characterization of the two confinement bottlenecks, unified bounds for coordinate-induced mixing with and without a leading core, and an empirical analysis reporting geometry, magnitude, and adaptation together. Broader language-understanding and generation results establish that the constructions support useful learning. The resulting perspective complements rank and parameter count with a precise account of what a spectral constraint restricts, what a relaxation changes, and what a regularizer can certify.

% Inactive modular snapshot; edit the root neurips_2026.tex instead.
\section{PEFT in pretrained singular coordinates}
\label{sec:coordinates}

For clarity, the main derivations use a square adapted matrix $W_0\in\mathbb R^{d\times d}$, as in the RoBERTa attention projections studied here. Let
\begin{equation}
W_0=U\Sigma V^\top
=U_L\Sigma_L V_L^\top+U_T\Sigma_T V_T^\top
\label{eq:pretrained}
\end{equation}
be a complete SVD with singular values in descending order. A declared cutoff $1\le k<d$ selects the leading $k$ directions, and the remaining $r=d-k$ directions form the tail. The cutoff is not the numerical rank: $\Sigma_T$ generally contains nonzero pretrained singular values. The pretrained weight $W_0$, including this tail, remains frozen. We define the effective update by $\Delta=W_{\mathrm{eff}}-W_0$ and its coordinates by
\begin{equation}
C=U^\top\Delta V=
\begin{bmatrix}C_{LL}&C_{LT}\\C_{TL}&C_{TT}\end{bmatrix}.
\label{eq:coordinates}
\end{equation}
Here $C_{LL}$ modifies the leading block, $C_{TT}$ acts within the tail, and $C_{LT},C_{TL}$ couple the two coordinate groups. We use \emph{cross mixing} for the last two blocks and \emph{off-tail allocation} for all three blocks outside $TT$. These terms are not interchangeable.

Orthogonal invariance and the disjoint blocks give
\begin{equation}
\fro{\Delta}^2=\fro{C}^2=\sum_{a,b\in\{L,T\}}\fro{C_{ab}}^2.
\label{eq:energy}
\end{equation}
For a nonzero update define $p_{ab}=\fro{C_{ab}}^2/\fro{\Delta}^2$, $p_{\cross}=p_{LT}+p_{TL}$, and $p_{\off}=p_{LL}+p_{\cross}$. The experiments report the off-tail \emph{norm ratio}, $q_{\off}=\sqrt{p_{\off}}$, and the relative magnitude $\rho_F=\fro{\Delta}/\fro{W_0}$. A mean of norm ratios cannot be squared to obtain the mean energy fraction. Absolute norms, normalized placement, and the aggregation across matrices must therefore be stated separately.

\paragraph{A common reference, not a method-specific score.}
The same fixed pretrained bases can measure any effective update, including LoRA. Block energies are invariant to orthogonal basis changes within each fixed subspace, although a diagonal coefficient family depends on the chosen paired basis. A singular-value tie across the cutoff makes the partition nonunique; one declared reference must then be held fixed. For rectangular matrices, complementary left and right projectors give an exhaustive decomposition (Appendix~\ref{app:rectangular}). When adapter insertion changes the initial effective matrix, displacement from $W_0$ differs from displacement since initialization; neither should be silently substituted for the other.

\paragraph{Structural change versus functional change.}
The coordinates describe weights. Leading-subspace drift additionally compares the top-ranked singular subspaces of $W_0$ and $W_0+\Delta$. Even a perfectly confined update can change that ranking if adapted tail singular values cross the boundary. Functional retention is a further question requiring behavior-level measurements; it does not follow from an unchanged block alone.

% Inactive modular snapshot; edit the root neurips_2026.tex instead.
\section{Confinement and flexibility}
\label{sec:confinement}

% Inactive modular snapshot; edit the root neurips_2026.tex instead.
\begin{figure}[t]
\centering
%
\begin{tabular}{cc}
\textbf{Fixed tail} & \textbf{Flexible within-tail core}\\[2pt]
$\begin{pmatrix}0&0\\0&\diag(h)\end{pmatrix}$ &
$\begin{pmatrix}0&0\\0&H\end{pmatrix}$\\[5pt]
Paired coefficients only & Same span, more internal freedom\\[9pt]
\textbf{Coordinate-relaxed core} & \textbf{Zero mixing, practical core}\\[2pt]
$\begin{pmatrix}C_{LL}&C_{LT}\\C_{TL}&C_{TT}\end{pmatrix}$ &
$\begin{pmatrix}C_{LL}&0\\0&C_{TT}\end{pmatrix}$\\[5pt]
Ambient scaling permits interaction & Leading modification may remain
\end{tabular}
\caption{Conceptual block structure in pretrained coordinates. Rotations change expressivity inside the tail, not its span. Ambient scaling can introduce off-tail entries. Suppressing both mixing operators restores tail-only support when the leading core is zero, but only block decoupling for the practical leading-plus-tail core. Nonzero blocks denote permitted locations, not arbitrary independent expressivity or measured data.}
\label{fig:blockschematic}
\end{figure}

We first remove ambient coordinate scaling and study updates supported entirely in the selected tail:
\begin{equation}
\Delta=U_T H V_T^\top.
\label{eq:tailcore}
\end{equation}
This family preserves the leading pretrained block and has zero cross mixing. Its available span, however, does not specify how much freedom is available \emph{inside} that span.

\paragraph{Fixed-tail coefficients.}
The fixed-direction construction sets $H=\diag(h)$, with $h\in\mathbb R^r$ trainable. We use LinLoRA as its implementation label, not as a claim to originate singular-value or fixed-basis coefficient tuning: these ideas appear in SVDiff and SVFT~\citep{han2023svdiff,lingam2024svft}. This construction changes paired singular-vector coefficients but cannot express off-diagonal interactions within the tail. The coefficients describe an additive update, not the entire singular spectrum of the adapted matrix.

\paragraph{Rotatable tail.}
The within-tail construction (OrthoLoRA) takes
\begin{equation}
H=R_U\diag(h)R_V^\top,\qquad R_U,R_V\in\mathbb O(r).
\label{eq:rotations}
\end{equation}
With unrestricted coefficients and full-tail rotations, the inner SVD makes every real $r\times r$ core representable. Rotations do not move the column spaces of $U_T$ and $V_T$. Practical partial rotations cover only $k_{\mathrm{vec}}$ of the $k_{\mathrm{val}}$ adapted directions, so their representable family is smaller than this full-core reference. Increasing flexibility also changes trainable capacity; this is not an equal-parameter comparison.

\begin{proposition}[Two distinct approximation bottlenecks]
\label{prop:approximation}
For an arbitrary target matrix $G\in\mathbb R^{d\times d}$, let $G_{TT}=U_T^\top G V_T$. Then
\begin{align}
\min_h\fro{G-U_T\diag(h)V_T^\top}^2
&=\fro{G}^2-\sum_{i=1}^r |(G_{TT})_{ii}|^2,\label{eq:fixederror}\\
\min_H\fro{G-U_T H V_T^\top}^2
&=\fro{G}^2-\fro{G_{TT}}^2.\label{eq:fullerror}
\end{align}
The improvement in the best attainable squared error equals the off-diagonal energy of $G_{TT}$. Both families retain the same unavoidable outside-tail residual.
\end{proposition}

The proof is an orthogonal-projection argument (Appendix~\ref{app:approximation}). This separates a \emph{within-tail coefficient bottleneck} from an \emph{outside-tail subspace bottleneck}. Rotations can address the former, not the latter. For example, a target $u_i v_j^\top$ with distinct tail indices is missed by paired diagonal coefficients but is representable by the full tail core. A leading-only target remains outside both families. These are exact approximation statements for a specified $G$, not assertions that a neural task optimum lies in the tail or that optimization reaches the best approximation.

\paragraph{What strict confinement preserves.}
For any fixed input $x$, $U_L^\top(W_0+\Delta)x=U_L^\top W_0x$ under~\eqref{eq:tailcore}. This is a local projection identity. Inputs to later layers can change during network adaptation, so it does not prove that pretrained behavior is preserved. The original leading singular pairs remain decoupled and are guaranteed to remain the top-$k$ pairs whenever
\begin{equation}
\sigma_k(W_0)>\op{\Sigma_T+H}.
\label{eq:nocrossing}
\end{equation}
Appendix~\ref{app:crossing} gives a tail-only counterexample when this separation fails.

% Inactive modular snapshot; edit the root neurips_2026.tex instead.
\section{Relaxation and leading--tail interaction}
\label{sec:mixing}

Trainable diagonal matrices $E,D$ act in the ambient output and input coordinates. For the tail-only family they produce $\Delta=E U_T H V_T^\top D$, relaxing strict confinement. We call this \emph{coordinate-relaxed adaptation}; the UI implementation labels are retained for reproducibility, without assuming a new invariance theorem for these learned scalers.

The practical construction also includes a fixed leading core. Both cases fit
\begin{equation}
\Delta=E\bigl(U_L S_L V_L^\top+U_T H V_T^\top\bigr)D,
\quad
S_L=\begin{cases}0&\text{tail-only reference},\\I&\text{practical leading-plus-tail form}.
\end{cases}
\label{eq:unified}
\end{equation}
The entire expression is an addition to $W_0$, including the leading term. Although $S_L$ is fixed, learning $E,D$ changes this contribution. Analyzing both cores exposes the distinction between suppressing cross-subspace interaction and eliminating leading adaptation.

\subsection{Mixing operators and a unified block bound}
Let $A=U^\top E U$ and $B=V^\top D V$, partitioned as in~\eqref{eq:coordinates}. The scalers are real diagonal, so $A,B$ are symmetric. Define
\begin{equation}
M_E=A_{LT}=U_L^\top E U_T,\quad
M_D=B_{LT}=V_L^\top D V_T,\quad
\mu=\op{M_E},\quad\nu=\op{M_D}.
\label{eq:mixingoperators}
\end{equation}
These are the $\mu_E,\nu_D$ quantities in the experimental record. Full within-tail orthogonal changes of basis do not change their operator or Frobenius norms. The effective blocks obey
\begin{equation}
C_{ab}=A_{aL}S_LB_{Lb}+A_{aT}HB_{Tb},\qquad a,b\in\{L,T\}.
\label{eq:blockidentity}
\end{equation}

\begin{proposition}[Interaction bounds with and without a leading core]
\label{prop:mixing}
Write $e=\op{E}$, $d_s=\op{D}$, $s=\op{S_L}$, and $h_s=\op{H}$. For~\eqref{eq:unified},
\begin{align}
\op{C_{LL}}&\le e s d_s+\mu h_s\nu, &
\op{C_{LT}}&\le e s\nu+\mu h_s d_s,\label{eq:blockbound1}\\
\op{C_{TL}}&\le\mu s d_s+e h_s\nu, &
\op{C_{TT}}&\le\mu s\nu+e h_s d_s.\label{eq:blockbound2}
\end{align}
If $\mu=\nu=0$, the cross blocks vanish. In addition, the leading block vanishes when $S_L=0$.
\end{proposition}

The proof applies submultiplicativity to~\eqref{eq:blockidentity} (Appendix~\ref{app:mixingproof}). The additional leading-core terms are the difference between the idealized and practical analyses. With $S_L=0$, leading--leading leakage requires mixing on both sides, while each cross block can be first order in one mixing quantity. With $S_L\ne0$, direct leading adaptation remains even when the cross terms vanish.

\begin{corollary}[Tail confinement and block decoupling]
\label{cor:confinement}
For $S_L=0$, define $C_{\mathrm{tail}}=\diag(0,C_{TT})$. Then
\begin{equation}
\op{C-C_{\mathrm{tail}}}
\le h_s\sqrt{(\mu\nu)^2+(\mu d_s)^2+(e\nu)^2}.
\label{eq:offbound}
\end{equation}
Consequently, $\mu,\nu\to0$ implies absolute tail confinement if $e,d_s,h_s$ stay bounded. If $e,d_s,h_s,s$ stay bounded for a general leading core, the same limit instead makes $C$ block-diagonal; it need not make $C_{LL}$ small.
\end{corollary}

For example, $E=D=I$, $S_L=I$, and $H=0$ gives zero mixing but $C_{LL}=I$. The practical update can therefore change the leading singular values or their ranking without cross mixing. Tail confinement, block decoupling, and ranked-subspace stability are three different conclusions.

\subsection{Regularization, scale, and effective geometry}
The soft intervention uses
\begin{equation}
\mathcal L=\mathcal L_{\mathrm{task}}
+\sum_{\ell\in\mathcal M}\left(\lambda_E\fro{M_{E,\ell}}^2+\lambda_D\fro{M_{D,\ell}}^2\right).
\label{eq:regularizer}
\end{equation}
Here $\mathcal M$ is the set of adapted matrices; the implementation sums, rather than averages, their penalties. Since $\op{M}\le\fro{M}$, each squared Frobenius penalty upper-bounds its squared mixing magnitude. The bounds are conditional statements about attained factors, not convergence guarantees for this objective. They control absolute blocks, not automatically the normalized fractions in Section~\ref{sec:coordinates}.

The factorization introduces a further distinction. Replacing $E$ by $aE$ and $D$ by $D/a$ leaves either effective update in~\eqref{eq:unified} unchanged, while replacing $(\mu,\nu)$ by $(a\mu,\nu/a)$. Thus a decrease in one raw mixing magnitude and an increase in the other is insufficient to establish redistribution of the matrix update. For the freely scalable tail-only core, $E,D\mapsto aE,aD$ and $H\mapsto H/a^2$ can reduce both raw penalties without changing $\Delta$. This latter limit violates the bounded-core condition when $a\to0$. These identities motivate measuring factor norms and effective blocks alongside the regularizer, and comparing its effect against update-magnitude control.

% Inactive modular snapshot; edit the root neurips_2026.tex instead.
\section{Controlling spectral interaction in practice}
\label{sec:experiments}

We organize the empirical analysis around three questions: how two-sided regularization changes effective update geometry, what differs under one-sided control, and whether the parameterizations support useful downstream adaptation. The mixing intervention is the primary experiment; the broader benchmarks provide supporting context.

\subsection{Study design and measurements}
The intervention uses RoBERTa-base on RTE, MRPC, CoLA, STS-B, and SST-2. Adapters modify the query, key, value, and attention-output projections in each of its 12 layers, giving 48 adapted matrices. The configuration adapts 256 coefficients with $k_{\mathrm{vec}}=0$, holding rotations out of the comparison. The analysis uses the practical leading-plus-tail form~\eqref{eq:unified}. The three conditions are A: no penalty, B: $\lambda_E=10^{-3},\lambda_D=0$, and C: $\lambda_E=\lambda_D=10^{-3}$. The backbone is frozen and the task head trains jointly with the adapter.

The analysis comprises seeds 42 and 17 for RTE, MRPC, CoLA, and STS-B, and seed 42 for SST-2: nine paired task/seed comparisons. Spectral quantities are measured at each run's validation-selected checkpoint, averaged over its 48 matrices, then averaged over seeds within a task. Runs, rather than layers, are the units of replication. Checkpoint selection is shared as a rule, not necessarily as an optimizer step. Appendix~\ref{app:details} gives the training settings, metric definitions, and complete three-condition results.

We report the relative effective update norm $\rho_F$, off-tail norm ratio $q_{\off}$, and the operator-norm $\sin\Theta$ distance between the leading pretrained and adapted singular subspaces. The latter is the sine of the \emph{largest} principal angle: a value near one means at least one direction is far apart, not that the whole leading subspace is orthogonal. The diagnostic cutoff is 512 leading directions and a 256-dimensional tail for each $768\times768$ matrix.

\begin{table}[t]
\centering
\caption{Two-sided mixing control for the practical spectral adapter. Each pair is A (unregularized) $\to$ C (two-sided penalty). Scores are multiplied by 100; $n$ is the number of seeds. Spectral entries are means over modules, then seeds, at validation-selected checkpoints. $q_{\off}$ is a norm ratio, not an energy fraction.}
\label{tab:mixingmain}
\begin{adjustbox}{max width=\linewidth}
\begin{tabular}{lccccc}
\toprule
Task & $n$ & Task score & $\rho_F$ & $q_{\off}$ & Left drift \\
\midrule
\tablerows{tables/mixing_summary_rows.tex}
\bottomrule
\end{tabular}
\end{adjustbox}
\end{table}

\subsection{Two-sided control changes both placement and magnitude}
Table~\ref{tab:mixingmain} shows a consistent shift toward lower off-tail allocation and subspace drift under the two-sided penalty. The reductions in $q_{\off}$, $\rho_F$, and both left and right drift occur in all nine paired task/seed comparisons. Task-level mean scores decrease by roughly $0.2$--$1.0$ points. Individual score differences are not uniformly negative across seeds; the experiment supports a descriptive trade-off, not statistical equivalence of task performance.

The geometric change is accompanied by pronounced shrinkage. Mean $\rho_F$ falls from approximately $0.28$--$0.38$ to $0.03$--$0.07$. The update remains nonzero, but that alone does not establish that the backbone update is necessary for the task score, since the head also trains. Meanwhile, the normalized off-tail ratio falls from approximately $0.95$ to $0.81$--$0.88$. Thus the intervention changes placement as well as magnitude relative to unregularized training, but does not achieve near-exact tail confinement.

\subsection{Which blocks change under mixing control?}
The aggregate off-tail ratio combines direct leading adaptation and cross interaction. To separate them, we reconstruct the four energy fractions from the saved per-matrix diagnostics, square before averaging, and retain equal weighting of modules within each run. Figure~\ref{fig:allocation} shows the resulting leading, cross, and tail allocation. Appendix~\ref{app:blockallocation} gives all four blocks and accounts for the small finite-precision difference between the recorded coordinate and ambient norms.

\begin{figure}[t]
\centering
% BEGIN GENERATED BLOCK FIGURE
% END GENERATED BLOCK FIGURE
\par\smallskip
{\small\textcolor{black!80}{\rule{1em}{1ex}} Leading--leading\quad
\textcolor{black!45}{\rule{1em}{1ex}} Cross\quad
\textcolor{black!15}{\rule{1em}{1ex}} Tail--tail}
\caption{Energy allocation in pretrained coordinates. Each bar sums to 100\% and averages per-matrix fractions, then seeds. A: no penalty; B: left-only; C: two-sided. Two-sided control reduces cross energy, but does not remove the leading block. Fractions describe relative placement, not absolute update size; Table~\ref{tab:mixingmain} reports the accompanying shrinkage.}
\label{fig:allocation}
\end{figure}

Without regularization, cross blocks contain approximately 43\% of update energy, the leading block 46\%, and the tail 10\%. Two-sided control reduces the cross share to 10--27\% and raises the tail share to 22--34\%. Nevertheless, the leading share rises in four of five task means and remains 44--64\% under C. This is the empirical counterpart of Proposition~\ref{prop:mixing}: suppressing cross interaction in a leading-plus-tail construction does not imply eliminating leading adaptation. An increased \emph{fraction} must not be read as increased absolute leading energy when the total update shrinks.

The decrease in cross share holds for all nine run-level pairs, but is not uniform across matrices: it occurs in 410 of 432 matched module/task/seed pairs. The off-tail norm ratio decreases in 332 pairs, and left and right drift in 260 and 256, respectively. These counts describe heterogeneity within trained models, not additional independent replicates. In particular, the consistent run means do not justify claiming that every layer's leading subspace is stabilized.

\subsection{One-sided control and factor asymmetry}
Across all nine paired comparisons, condition B decreases $\mu_E$ and increases $\nu_D$ relative to A. Its task score is highest among the three task-level means on RTE, MRPC, and CoLA, while substantial off-tail allocation and drift remain. This pattern is consistent with the distinction between constraining one mixing operator and controlling the complete update. It does not establish that learning moves through the unpenalized side: reciprocal scaler changes can produce the same factor pattern without changing $\Delta$ (Section~\ref{sec:mixing}). Appendix~\ref{app:mixingtable} reports both factors alongside the effective-matrix diagnostics.

\paragraph{Interpreting normalized geometry.}
The changed norm ratios exclude an explanation based only on multiplying the \emph{same trained matrix} by a scalar. They do not exclude the possibility that magnitude-only regularization during training would produce similar geometry. The distinction is between an observed change in normalized placement and a magnitude-independent benefit, which requires a norm-matched training control. These selected endpoints also concern update geometry, not the time at which mixing arises during learning.

The size of the blocks also matters. An isotropically oriented perturbation has expected leading, cross, and tail energy fractions $k^2/d^2$, $2kr/d^2$, and $r^2/d^2$. At this cutoff these are $4/9$, $4/9$, and $1/9$. The unregularized allocation is close to this dimension-only reference; a large aggregate off-tail fraction alone therefore does not demonstrate preferential use of leading directions. The reference is neither a fitted model nor a significance test (Appendix~\ref{app:dimension}).

\subsection{Downstream viability of the instruments}
Table~\ref{tab:glue_results_base_large} presents the full six-task GLUE evaluation on both RoBERTa backbones. Its role is to test whether the instruments support useful adaptation, not to rank new adapters. The shared five-task average is 84.20 versus 83.62 for the rotated and diagonal cores on RoBERTa-base, and 87.36 versus 87.54 on RoBERTa-large. The ordering reverses: rotations are not uniformly better. This is compatible with Proposition~\ref{prop:approximation}, which characterizes available directions rather than guaranteeing a benefit from using them. The mixing experiment disables rotations, so its measured regularization effect does not require a rotated core.

% Inactive modular snapshot; edit the root neurips_2026.tex instead.
\begin{table}[t]
\centering
\caption{RoBERTa validation results on all six evaluated GLUE tasks. Entries are reported mean$\pm$standard deviation, multiplied by 100. VeRA, LoRA, and RandLoRA are imported comparisons~\citep{albert2025randlora}; UI rows use the spectral protocol. $\dagger$MRPC is accuracy for imported rows and F1 for UI rows, so it is excluded from Avg5 for \emph{every} method. Avg5 averages the other five displayed means, not the full GLUE benchmark.}
\label{tab:glue_results_base_large}
\begin{adjustbox}{max width=\linewidth}
\begin{tabular}{@{}lccccccc@{}}
\toprule
Method & SST-2 & MRPC$^\dagger$ & CoLA & QNLI & RTE & STS-B & Avg5 \\
\midrule
% BEGIN GENERATED GLUE ROWS
% END GENERATED GLUE ROWS
\bottomrule
\end{tabular}
\end{adjustbox}
\end{table}


The comparison includes imported VeRA, LoRA, and RandLoRA results~\citep{albert2025randlora}, not uniformly retuned baselines. Their MRPC scores are accuracy, whereas the spectral protocol uses F1; we retain the separate results but exclude MRPC from every average. Even the shared-metric average is not a full GLUE score or a matched comparison: MNLI and QQP are absent, tuning protocols differ, and rotations change capacity. The table does not establish superiority over the wider spectral-PEFT literature or an isolated rotation benefit.

GPT-2 Medium generation results provide a check beyond classification and regression. An exploratory Gemma/Llama study measures adaptation and retention along a separate axis. Neither connects matched spectral measurements to its functional outcomes; these evaluations therefore do not establish that the RoBERTa mixing mechanism transfers to modern generative models.

\paragraph{Computational scope.}
Compact trainable capacity does not imply a compact spectral reference. For a square $d\times d$ layer, dense SVD initialization costs $O(d^3)$ and the practical implementation stores $2d^2$ frozen basis entries. Its unmerged adapter uses full-basis products with $O(bd^2)$ cost for $b$ token vectors, in addition to the frozen base projection. The merged model has the same linear-layer inference structure as the base model. Appendix~\ref{app:cost} separates setup, storage, trainable capacity, forward computation, regularization, and merging costs; we do not infer training speed or total memory savings from parameter counts alone.

% Inactive modular snapshot; edit the root neurips_2026.tex instead.
\section{Related work}
\label{sec:related}

\paragraph{Efficient parameterizations.}
LoRA~\citep{hu2022lora} restricts additive updates through low-rank factors. AdaLoRA adaptively allocates an update's singular-value budget across matrices~\citep{zhang2023adalora}; its SVD-like update factors should not be confused with a fixed SVD of the pretrained weight. VeRA, DoRA, FourierFT, and RandLoRA explore alternative trainable structures and efficiency trade-offs~\citep{kopiczko2023vera,liu2024dora,gao2024parameter,albert2025randlora}. Intrinsic-dimension studies likewise show that learning can succeed in restricted parameter spaces~\citep{li2018measuring,aghajanyan2020intrinsic}. Low dimension, however, does not determine placement relative to the pretrained spectrum.

\paragraph{Pretrained spectral information.}
SVDiff tunes pretrained singular values, while SVFT learns sparse coefficients of pretrained singular-vector outer products~\citep{han2023svdiff,lingam2024svft}. These directly precede the fixed-basis constructions studied here. PiSSA initializes trainable factors from leading components, whereas MiLoRA initializes them from minor components~\citep{meng2024pissa,wang2025milora}. Such initialization does not, by itself, constrain the learned factors to remain in their initial spans. LoRA-XS and CorDA provide further decomposition-guided parameterizations~\citep{balazy2024lora,yang2024corda}. We use the pretrained frame to distinguish these choices rather than claim that this frame or singular-coefficient tuning is new.

\paragraph{Orthonormality versus confinement.}
SORSA is closely related: it trains principal singular values and vectors, freezes a residual weight, and regularizes the trainable vectors' Gram matrices toward identity~\citep{cao2025sorsa}. It also analyzes spectral changes. Our within-tail rotations instead preserve prescribed pretrained spans exactly, while our mixing penalty acts on overlaps induced by ambient scalers relative to those spans. Orthonormal vectors can span a different subspace, so a Gram penalty and a pretrained leading--tail mixing penalty constrain different objects. This distinction defines our analytical focus; it is not evidence of better downstream performance than SORSA, which is not included in the present empirical comparison.

\paragraph{Geometry, intervention, and retention.}
\citet{shuttleworth2025illusion} study spectral differences between LoRA and full fine-tuning, including interventions on intruder directions. A recent preprint additionally studies spectral thresholds and norm-matched interventions on selected layers~\citep{xie2026intruder}. CLoRA uses subspace regularization to control adaptation in continued training~\citep{lu2025controlled}; \citet{biderman2024lora} study the adaptation--forgetting trade-off. Our analysis complements these perspectives by resolving a specific source of interaction---ambient scaling around a pretrained spectral core---and separating its factor-level bounds from effective-update measurements. Neither a new singular direction nor a changed coordinate block has a universal functional interpretation without task-level evidence.

% Inactive modular snapshot; edit the root neurips_2026.tex instead.
\section{Discussion and limitations}
\label{sec:discussion}

\paragraph{What the tail reference does and does not imply.}
Leading and tail directions are defined by singular-value order, not by their semantic importance. The projection identities hold for any target perturbation and either selected spectral region; they do not require an alignment assumption. However, the amount of task-relevant signal captured by a region is an empirical question. Without equal-dimensional leading, tail, and mixed/random controls, the chosen tail cannot be identified as uniquely suitable for adaptation. Likewise, increased within-tail expressivity need not improve a trained model: the useful directions, optimization, and statistical cost of added capacity all matter.

\paragraph{Scope of the intervention.}
The five-task intervention has limited seed replication, validation-selected endpoints, and no magnitude-matched training or head-only control. It establishes that the tested penalty changes geometry together with magnitude, but not that its spectral effect is necessary for adaptation or independent of shrinkage. Factor penalties also admit rescaling ambiguities, so the effective matrix and the factor norms in the bounds remain essential measurements. The approximation results concern representability, not optimization or generalization; the mixing bounds are conditional on attained factors, and ranked-subspace stability additionally requires spectral separation.

\paragraph{External validity and comparison coverage.}
The mechanistic measurements use RoBERTa-base, not contemporary large-model instruction tuning or reasoning. The broader benchmark comparisons are not a uniformly tuned evaluation against PiSSA, MiLoRA, AdaLoRA, DoRA, and SORSA. The generation and retention results show additional applications but do not supply that missing geometric comparison; the larger-model study also lacks a scaler-only UILinLoRA counterpart. These limitations preclude claims of universal tail preference, a consistent advantage from rotations, or functional preservation induced by mixing control. They motivate transferring the same controlled contrasts to a generative backbone, rather than treating an additional leaderboard score as evidence for the mechanism.

\section{Conclusion}
\label{sec:conclusion}
We studied spectral fine-tuning through the location and structure of the effective update in pretrained singular coordinates. A diagonal tail core and a flexible tail core share an outside-tail bottleneck but differ in their within-tail expressivity. Ambient scaling introduces interaction whose control depends on the underlying core: vanishing mixing gives tail confinement for a tail-only update, but only cross-block decoupling when direct leading adaptation is present. The empirical intervention reduces cross energy while the leading block retains a substantial share, alongside update shrinkage and a modest task-score cost. Understanding spectral adaptation therefore requires distinguishing available directions, cross-subspace interaction, and effective magnitude. These distinctions make the parameterizations useful instruments for studying PEFT, independently of which adapter scores highest.

% Inactive modular snapshot; edit the root neurips_2026.tex instead.
\section*{Reproducibility statement}
Complete proofs appear in Appendix~\ref{app:proofs}, training and metric definitions in Appendix~\ref{app:details}, and computational costs in Appendix~\ref{app:cost}. Equation~\eqref{eq:unified} specifies the practical parameterization separately from the tail-only reference. The accompanying per-task/seed summary records and table-generation script reproduce the numerical aggregation for the mixing study; a separate synthetic-matrix script checks the algebraic identities and boundary cases. These checks do not constitute a rerun of model training. Supporting benchmark tables distinguish imported comparisons from the spectral-variant results.

\section*{AI use statement}
Generative AI tools were used to assist with restructuring and drafting the manuscript, literature lookup, mathematical formulation and proof drafting, interpretation of existing experimental records, and design of additional experiments. AI assistance was also used to write scripts for consistency checks and table reconstruction. The numerical experimental results reported here come from existing research records, not generated or simulated outcomes attributed to unperformed training runs.

\begin{samepage}
\section*{Ethics statement}
The study analyzes adaptation of existing language models using existing benchmarks. Weight-space stability does not certify factual correctness, fairness, privacy, or safe behavior. Efficient adaptation can reduce training costs but can also facilitate harmful specialization or preserve biases present in a base model. Our structural claims should not be interpreted as deployment guarantees.\par
\end{samepage}

\bibliography{reference}
\bibliographystyle{iclr2027_conference}
\clearpage
\appendix
% Inactive modular snapshot; edit the root neurips_2026.tex instead.
\section{Proofs and boundary cases}
\label{app:proofs}

\subsection{Fixed coefficients and within-tail flexibility}
\label{app:approximation}
Let $\widehat G=U^\top G V$ and partition it into the four coordinate blocks. Orthogonal invariance gives
\begin{equation}
\fro{G-U_T H V_T^\top}^2
=\fro{\widehat G_{LL}}^2+\fro{\widehat G_{LT}}^2
+\fro{\widehat G_{TL}}^2+\fro{\widehat G_{TT}-H}^2.
\end{equation}
For a free core the last term is minimized at $H=\widehat G_{TT}$. For a diagonal core it is minimized by $h_i=(\widehat G_{TT})_{ii}$, leaving the off-diagonal entries unchanged. Substitution proves Proposition~\ref{prop:approximation}. An SVD of the minimizing free core gives the representation~\eqref{eq:rotations}; this is an existence statement, not a guarantee for a particular optimizer or partial-rotation block.

If the free core is additionally constrained to rank at most $q\le r$, applying truncated SVD to $\widehat G_{TT}$ gives minimum squared error
\begin{equation}
\fro{G}^2-\sum_{j=1}^q\sigma_j(\widehat G_{TT})^2.
\end{equation}
This expression keeps the outside-tail residual explicit. It should not be replaced by the best unrestricted rank-$q$ approximation error of the entire $G$, whose singular directions may lie outside the selected tail.

\subsection{Proof of the unified mixing bounds}
\label{app:mixingproof}
Because $U,V$ are orthogonal,
\begin{equation}
C=A\begin{bmatrix}S_L&0\\0&H\end{bmatrix}B.
\end{equation}
Direct multiplication gives
\begin{align}
C_{LL}&=A_{LL}S_LB_{LL}+A_{LT}HB_{TL},\\
C_{LT}&=A_{LL}S_LB_{LT}+A_{LT}HB_{TT},\\
C_{TL}&=A_{TL}S_LB_{LL}+A_{TT}HB_{TL},\\
C_{TT}&=A_{TL}S_LB_{LT}+A_{TT}HB_{TT}.
\end{align}
Symmetry yields $\op{A_{TL}}=\op{A_{LT}}=\mu$ and $\op{B_{TL}}=\op{B_{LT}}=\nu$. Orthogonal compression bounds the diagonal blocks of $A$ by $e$ and those of $B$ by $d_s$. Apply the triangle inequality and submultiplicativity to each product to obtain~\eqref{eq:blockbound1}--\eqref{eq:blockbound2}.

For the tail-only core, $s=0$ and
\begin{equation}
C-C_{\mathrm{tail}}=\begin{bmatrix}C_{LL}&C_{LT}\\C_{TL}&0\end{bmatrix}.
\end{equation}
For any block vector $(x,y)$, the norms of the two output blocks are bounded componentwise by
\begin{equation}
\begin{bmatrix}\op{C_{LL}}&\op{C_{LT}}\\\op{C_{TL}}&0\end{bmatrix}
\begin{bmatrix}\|x\|_2\\\|y\|_2\end{bmatrix}.
\end{equation}
Taking a supremum and bounding the norm of this nonnegative $2\times2$ matrix by its Frobenius norm proves
\begin{equation}
\op{C-C_{\mathrm{tail}}}
\le\sqrt{\op{C_{LL}}^2+\op{C_{LT}}^2+\op{C_{TL}}^2}.
\end{equation}
Substitution gives~\eqref{eq:offbound}. If the stated factor norms remain bounded, its right-hand side goes to zero as $\mu,\nu\to0$.

For the general core, let $C_{\mathrm{diag}}=\diag(C_{LL},C_{TT})$. The off-diagonal block matrix satisfies the exact identity
\begin{equation}
\op{C-C_{\mathrm{diag}}}=\max\{\op{C_{LT}},\op{C_{TL}}\}.
\end{equation}
The cross-block bounds therefore imply convergence to the set of block-diagonal updates when $e,d_s,h_s,s$ are bounded and $\mu,\nu\to0$. The term $A_{LL}S_LB_{LL}$ need not vanish. This proves Corollary~\ref{cor:confinement} and states precisely the different limits in the two cases.

\subsection{Confinement does not prevent singular-value crossings}
\label{app:crossing}
For a tail-only update without ambient scaling,
\begin{equation}
U^\top(W_0+\Delta)V=
\begin{bmatrix}\Sigma_L&0\\0&\Sigma_T+H\end{bmatrix}.
\end{equation}
The singular values are the union of those of the two diagonal blocks, which proves the sufficient condition~\eqref{eq:nocrossing}. Consider $W_0=\diag(2,1)$, cutoff $k=1$, and $\Delta=\diag(0,2)$. The update is strictly tail-supported with no cross mixing and no leading-block modification, but $W_0+\Delta=\diag(2,3)$ has the other coordinate as its top singular direction. Its largest-principal-angle sine relative to the original leading direction is one. Thus block decoupling and stability of the top-ranked subspace are distinct even in a two-dimensional example.

\subsection{Rectangular weights and complete energy accounting}
\label{app:rectangular}
For $W_0\in\mathbb R^{m\times n}$, select $k\le\min(m,n)$ leading singular vectors and define
\begin{equation}
P_L=U_LU_L^\top,\quad P_T=I_m-P_L,\qquad
Q_L=V_LV_L^\top,\quad Q_T=I_n-Q_L.
\end{equation}
Then $\Delta_{ab}=P_a\Delta Q_b$ gives
\begin{equation}
\Delta=\sum_{a,b}\Delta_{ab},\qquad
\fro{\Delta}^2=\sum_{a,b}\fro{\Delta_{ab}}^2.
\end{equation}
Pairwise Frobenius inner products vanish because complementary projectors multiply to zero. The tail here is the full complement and includes any unmatched null directions. A chosen adapter may span only a subset of it; that available subspace must not be confused with the diagnostic complement. Complete orthogonal bases recover the corresponding coordinate blocks. The square attention-matrix experiment avoids this issue, but an economy-SVD diagnostic is not automatically exhaustive for other layers.

\subsection{A dimension-only reference for allocation}
\label{app:dimension}
Let the vectorized nonzero update have a rotationally invariant direction in the $d^2$-dimensional matrix space. By symmetry each squared coordinate has expected fraction $1/d^2$ of the total. The leading, cross, and tail blocks have $k^2$, $2kr$, and $r^2$ coordinates, respectively. Their expected energy fractions are therefore $k^2/d^2$, $2kr/d^2$, and $r^2/d^2$. For $r/d=1/3$, these are $4/9$, $4/9$, and $1/9$, with off-tail expectation $8/9$. The square root $\sqrt{8/9}\simeq0.9428$ is a reference on the norm-ratio scale, not the expectation $\mathbb E[q_{\off}]$. These calculations are a null reference, not a fitted model or a statistical test of the observed updates.

\subsection{Factor-scale identities}
Write the inner matrix in~\eqref{eq:unified} as $K$. For any $a>0$, $(aE)K(D/a)=EKD$. Because each mixing operator is linear in its respective scaler, this transformation multiplies its two magnitudes by $a$ and $1/a$. For $S_L=0$, the additional identity
\begin{equation}
(aE)U_T(H/a^2)V_T^\top(aD)=E U_T H V_T^\top D
\end{equation}
shrinks both squared Frobenius penalties by $a^2$. It can make a raw factor penalty arbitrarily small with a fixed effective update if the core is allowed to diverge. It does not refute the bounded-factor confinement statement, nor does it imply that the optimizer follows this path. It motivates measuring the actual update and the norms appearing in the bound.

% Inactive modular snapshot; edit the root neurips_2026.tex instead.
\section{Experimental details}
\label{app:details}

\subsection{Backbones, tasks, and placement}
The broader language-understanding evaluation uses RoBERTa-base and RoBERTa-large on a six-task GLUE subset: CoLA, SST-2, MRPC, STS-B, QNLI, and RTE~\citep{wang2018glue}. The mixing intervention uses RoBERTa-base on five of these tasks, excluding QNLI. Neither evaluation covers MNLI or QQP. Our task metrics are Matthews correlation for CoLA, accuracy for SST-2/QNLI/RTE, F1 for MRPC, and Pearson correlation for STS-B. The imported RandLoRA-paper comparisons instead report MRPC accuracy, which is why Table~\ref{tab:glue_results_base_large} excludes MRPC from the shared average. Results use the official validation splits. Pretrained backbone parameters are frozen, and adapters and the prediction head train jointly. The adapted modules are the query, key, value, and attention-output projections; feed-forward projections remain frozen.

\subsection{Training configuration}
The training recipe uses AdamW with separate adapter and head learning rates, linear decay, a 6\% warmup fraction, batch size 64, and maximum sequence length 256. Table~\ref{tab:recipe} gives the task-specific settings recorded for RoBERTa-base. These task-specific configurations are not a common hyperparameter search across all methods in the supporting benchmark table.

\begin{table}[ht]
\centering
\caption{RoBERTa-base training settings from the experimental protocol. The scaler value initializes both ambient diagonal scalers.}
\label{tab:recipe}
\begin{tabular}{lccccc}
\toprule
Task & Epochs & Head LR & Adapter LR & Scaler & Coefficients\\
\midrule
CoLA & 80 & $5\cdot10^{-3}$ & $3\cdot10^{-2}$ & $10^{-2}$ & $10^{-1}$\\
SST-2 & 40 & $10^{-2}$ & $4\cdot10^{-2}$ & $10^{-1}$ & $10^{-1}$\\
MRPC & 30 & $10^{-3}$ & $5\cdot10^{-2}$ & $10^{-1}$ & $10^{-1}$\\
STS-B & 60 & $5\cdot10^{-3}$ & $10^{-2}$ & $10^{-1}$ & $10^{-1}$\\
QNLI & 25 & $10^{-3}$ & $2\cdot10^{-2}$ & $10^{-1}$ & $10^{-1}$\\
RTE & 90 & $5\cdot10^{-4}$ & $10^{-2}$ & $10^{-2}$ & $10^{-2}$\\
\bottomrule
\end{tabular}
\end{table}

Adapter dropout is zero and the nominal adapter scaling coefficient is one. The broad GLUE configurations use $k_{\mathrm{val}}=256$, with $k_{\mathrm{vec}}=0$ for UILinLoRA and 64 for UIOrthoLoRA. The mixing study uses 256 and zero, respectively. Compact rotations, when enabled, use an orthogonal parameterization; only the selected sub-block is rotated. With nonzero initialized scalers and coefficients, the practical adapter need not have zero initial delta. The measured displacement from $W_0$ includes that initial contribution.

\subsection{Seeds and checkpoint selection}
The broader spectral evaluation protocol reports six seeds $\{42,123,2021,17,31415,1054\}$; Table~\ref{tab:glue_results_base_large} retains its reported means and across-seed standard deviations. The imported comparison uses five runs. The mixing analysis uses the nine task/seed pairs with archived paired summary and module-level records: seeds 42 and 17 for RTE, MRPC, CoLA, and STS-B, and 42 for SST-2. These records permit independent reaggregation of the intervention, whereas the broad GLUE table is a report of the original evaluation, not a newly reconstructed six-seed experiment.

The ablation runner reloads the best validation checkpoint before computing final scores and spectral metrics. Consequently, the selected checkpoints can occur at different steps: the RTE seed-42 A/B/C records, for example, select steps 312, 3198, and 1755. The endpoint analysis does not use stepwise loss-log entries as measurements of spectral trajectories.

\subsection{Metric definitions and aggregation}
For each adapted matrix, the reported metrics are
\begin{align}
\rho_F&=\fro{\Delta}/\fro{W_0},&
q_{\off}&=\fro{C_{\off}}/\fro{C},\\
\mathrm{Leak}_{LL}&=\op{C_{LL}}/\op{\Delta},&
C_{\off}&=\begin{bmatrix}C_{LL}&C_{LT}\\C_{TL}&0\end{bmatrix}.
\end{align}
The mixing magnitudes are operator norms, while their training penalties use squared Frobenius norms. Leading left and right drift use the largest-principal-angle sine between the leading subspaces of $W_0$ and $W_0+\Delta$. The tables average each metric across the 48 adapted modules within a run, then across available seeds. They do not pool module energies before normalization. Because each model is a single trained replicate, 48 modules do not constitute 48 independent experimental seeds.

\subsection{Complete mixing results}
\label{app:mixingtable}
The accompanying \texttt{scripts/build\_mixing\_tables.py} generates Tables~\ref{tab:mixingmain} and~\ref{tab:mixingfull} from the included per-task/seed CSV summaries. These records retain checkpoint steps and metric columns. The analysis includes the paired A/B/C runs and excludes separate smoke tests and C-only reruns. This makes the aggregation independently checkable without rerunning training.
\begin{table}[ht]
\centering
\caption{Complete mixing-ablation means. A: no penalty; B: left-only penalty; C: two-sided penalty. $n$ is the seed count. Scores are in native metric units. All spectral values are module means followed by seed means.}
\label{tab:mixingfull}
\begin{adjustbox}{max width=\linewidth}
\begin{tabular}{llcccccccc}
\toprule
Task & Condition & $n$ & Score & $\mu_E$ & $\nu_D$ & $\rho_F$ & $q_{\off}$ & Drift$_U$ & Drift$_V$\\
\midrule
\tablerows{tables/mixing_full_rows.tex}
\bottomrule
\end{tabular}
\end{adjustbox}
\end{table}

\subsection{Reconstructing four-block allocation}
\label{app:blockallocation}
The 27 run summaries are accompanied by 1,296 module-level records: 48 adapted matrices for each A/B/C run. Their saved Frobenius leakage ratios use $\|\Delta\|_F$ in the denominator, whereas the off-tail ratio uses $\|C\|_F$. Write these recorded quantities as $\ell_{ab}$ for $ab\in\{LL,LT,TL\}$ and $q$ for the off-tail ratio. To obtain an allocation normalized in the recorded coordinate frame, we compute per matrix
\begin{equation}
\widehat p_{ab}=q^2\frac{\ell_{ab}^2}{\ell_{LL}^2+\ell_{LT}^2+\ell_{TL}^2},\qquad
\widehat p_{TT}=1-q^2.
\label{eq:recordedallocation}
\end{equation}
For exact orthogonal bases these agree with Section~\ref{sec:coordinates}. With saved finite-precision bases, the inferred discrepancy $\left|(\ell_{LL}^2+\ell_{LT}^2+\ell_{TL}^2)/q^2-1\right|$ is at most $2.64\cdot10^{-4}$ across the archive. The normalization makes the four recorded fractions sum to one without assuming perfect numerical orthogonality. Diagnostic denominators also include a $10^{-12}$ numerical stabilizer. We square each matrix's ratios before averaging, never a task-mean ratio.

\begin{table}[ht]
\centering
\caption{Four-block energy percentages underlying Figure~\ref{fig:allocation}, averaged over modules and then seeds. $n$ is the number of seeds. The two middle columns sum to the cross share. Each row sums to 100\% up to rounding.}
\label{tab:blockallocation}
\begin{tabular}{llccccc}
\toprule
Task & Condition & $n$ & Leading--leading & Leading--tail & Tail--leading & Tail--tail\\
\midrule
% BEGIN GENERATED BLOCK TABLE
% END GENERATED BLOCK TABLE
\bottomrule
\end{tabular}
\end{table}

The analysis script reconciles every shared summary metric with the mean of its 48 saved module values before calculating these fractions. It excludes smoke tests and unmatched C-only reruns, as in the summary analysis. The 432 matched A/C module pairs measure within-run heterogeneity; replication remains nine task/seed pairs, not hundreds of independent training runs.

\subsection{Assets and compute reporting}
The experiments use public benchmarks, pretrained models, and the PyTorch/Hugging Face software ecosystem. No new dataset is introduced. The experimental record reports NVIDIA A100-SXM4-40GB and RTX-3090 hardware. The mixing summaries include per-run runtime and memory fields, but not a complete project-wide compute accounting or a controlled throughput comparison across methods. The following analysis therefore reports arithmetic and storage costs separately from empirical efficiency claims.

% Inactive modular snapshot; edit the root neurips_2026.tex instead.
\section{Computational cost of the spectral instruments}
\label{app:cost}
We distinguish trainable capacity from frozen storage, setup cost, and the computation required to apply the adapter. Let a square layer have width $d$, selected tail dimension $r$, and rotated sub-block dimension $q\le r$. Let $b$ count the token vectors processed by that layer in one forward call. All costs below are per adapted matrix and exclude the base model's own storage and forward computation, the task head, and hardware-dependent constants.

\paragraph{SVD initialization and frozen storage.}
A dense SVD of a $d\times d$ matrix has conventional $O(d^3)$ arithmetic cost. For an $m\times n$ matrix with $p=\min(m,n)$, the corresponding economy-SVD cost is $O(mnp)$ and its two bases contain $(m+n)p$ entries. This is a one-time setup cost per distinct weight decomposition, not a decomposition at every optimization step. The practical implementation computes the SVD in float32 before casting stored factors to the adapter dtype, so initialization workspace can exceed persistent basis storage.

A tail-only implementation needs $2dr$ basis entries to apply its update. The implemented leading-plus-tail construction retains the complete bases, totaling $2d^2$ entries, plus $O(d)$ fixed coefficients and any orthogonal-map state. Freezing these entries avoids gradients and optimizer moments for them; it does not remove their memory cost. Setting the leading coefficients to zero in a full-basis implementation does not itself discard those bases or shorten its matrix products. Diagnostics that retain an additional full reference must also be included in memory accounting.

\paragraph{Trainable entries versus effective degrees of freedom.}
Fixed-tail coefficients require $r$ trainable entries. Two dense $q\times q$ rotation parameters add $2q^2$ stored trainable entries; their orthogonal outputs have only $q(q-1)$ combined manifold degrees of freedom. Ambient input/output scalers add $2d$ entries. Thus the practical unrotated and partially rotated forms store $r+2d$ and $r+2d+2q^2$ trainable entries, respectively, before the task head. A directly trained unrestricted tail core instead has $r^2$ entries. These counts describe different parameterizations, not equal-capacity interventions. For example, $d=768,r=256$ gives 1,792 trainable entries per unrotated coordinate-scaled matrix but 1,179,648 frozen basis entries. Aggregate benchmark parameter counts in Appendix~\ref{app:downstream} are retained as reported, rather than substituted for a tensor-level memory inventory.

\paragraph{Forward computation.}
Applying a diagonal tail update through its thin factors costs $O(bdr)$; ambient scaling adds only $O(bd)$. If partial rotations are materialized into the basis first, the two basis products add $O(dq^2)$ per reconstruction, in addition to the orthogonal-map construction cost. Dense $q\times q$ maps can incur $O(q^3)$ arithmetic. A dense tail core applied between thin basis products instead costs $O(bdr+br^2)$.

The practical forward path concatenates the leading and tail factors and applies two full-basis products. Its adapter branch therefore costs $O(bd^2)$ even when $q=0$, plus rotation construction when enabled. This is additional to the $O(bd^2)$ base projection. A rank-$\ell$ LoRA branch costs $O(bd\ell)$, so a smaller trainable count need not mean a cheaper forward pass. Backpropagation also differentiates through the adapter operations, and end-to-end peak memory depends on retained activations, dtype, optimizer state, and checkpointing. The arithmetic counts alone do not establish a measured throughput ratio.

\paragraph{Regularization and diagnostics.}
Directly forming $M_E=U_L^\top E U_T$ and $M_D=V_L^\top D V_T$ costs $O(dkr)$ for a leading dimension $k=d-r$; evaluating their squared Frobenius penalties then costs $O(kr)$. These operations and their gradients are additional training work when the penalty is enabled. They do not require an SVD of the adapted weight. By contrast, exact ranked-subspace diagnostics use a new dense SVD of $W_0+\Delta$ and can cost $O(d^3)$ per measured layer/checkpoint. Diagnostic time should be reported separately from optimization time and should not be charged only to the spectral methods in a controlled comparison.

\paragraph{Merging and deployment.}
At evaluation with dropout disabled, the additive update can be materialized and merged into $W_0$. A thin diagonal-tail construction costs $O(d^2r)$ to materialize, whereas the current full-basis product costs $O(d^3)$, plus the rotation and scaler operations; dense materialization requires $O(d^2)$ output storage. After merging and removing the adapter branch, the linear layer has the base model's shape and inference-time matrix multiplication. This removes the extra branch at deployment, not the preceding SVD, training, or merge costs.

% Inactive modular snapshot; edit the root neurips_2026.tex instead.
\section{Supporting downstream evaluations}
\label{app:downstream}
\subsection{Generation}
\label{app:generation}

Table~\ref{tab:gpt2_generation_appendix} reports GPT-2 Medium results on E2E data-to-text generation~\citep{novikova2017e2e}. The fixed-core and partially rotated coordinate-scaled constructions attain BLEU scores of 68.7 and 68.8, respectively, with 0.078M and 0.35M reported trainable parameters. These results extend the practical setting beyond classification and regression, but contain no block-allocation measurements or mixing intervention. The non-UI comparisons are taken from \citet{gao2024parameter}, whose LoRA and FourierFT generation entries average three final-epoch runs. We retain source-reported uncertainties without treating them as a common inferential comparison. This is supporting task-level evidence, not a uniformly rerun geometry experiment.
% Inactive modular snapshot; edit the root neurips_2026.tex instead.
\begin{table}[ht]
\centering
\caption{Supporting GPT-2 Medium E2E generation results. Non-UI comparison values are from \citet{gao2024parameter}, including their imported full-fine-tuning and adapter results; $L$ and $H$ denote Lin-style and Houlsby-style adapters. Uncertainties retain their source reporting conventions. These task scores are not measurements of spectral mixing.}
\label{tab:gpt2_generation_appendix}
\begin{adjustbox}{max width=\linewidth}
\begin{tabular}{@{}l r c c c c c@{}}
\toprule
\textbf{Method} & \textbf{\# Params} & \textbf{BLEU} & \textbf{NIST} & \textbf{METEOR} & \textbf{ROUGE-L} & \textbf{CIDEr} \\
\midrule
Full FT          & 354.92M & 68.2           & 8.62           & 46.2           & 71.0           & 2.47 \\
Adapter$^L$      & 0.37M   & 66.3           & 8.41           & 45.0           & 69.8           & 2.40 \\
Adapter$^L$      & 11.09M  & 68.9           & 8.71           & 46.1           & 71.3           & 2.47 \\
Adapter$^H$      & 11.09M  & 67.3$_{\pm.6}$ & 8.5$_{\pm.07}$ & 46.0$_{\pm.2}$ & 70.7$_{\pm.2}$ & 2.44$_{\pm.01}$ \\
LoRA             & 0.35M   & 68.9$_{\pm.3}$  & 8.76$_{\pm.06}$ & 46.6$_{\pm.1}$ & 71.5$_{\pm.1}$ & 2.53$_{\pm.03}$ \\
FourierFT     & 0.048M  & 69.1$_{\pm.1}$ & 8.82$_{\pm.05}$ & 47.0$_{\pm.3}$ & 71.8$_{\pm.1}$ & 2.51$_{\pm.02}$ \\
UIOrthoLoRA   & 0.35M   & 68.8$_{\pm.1}$ & 8.729$_{\pm.06}$ & 46.71$_{\pm.3}$ & 71.6$_{\pm.02}$ & 2.47$_{\pm.01}$ \\
UILinLoRA     & 0.078M  & 68.7$_{\pm.1}$ & 8.732$_{\pm.1}$ & 46.52$_{\pm.3}$ & 71.45$_{\pm.01}$ & 2.46$_{\pm.02}$ \\
\bottomrule
\end{tabular}
\end{adjustbox}
\end{table}


\section{Exploratory adaptation and retention}
\label{app:retention}

The additional retention study measures task adaptation and behavior retention for instruction-tuned Gemma-3-12B and Llama-3.2-3B models~\citep{gemma3_2025,llama3_2024}. It remains exploratory support rather than evidence for the spectral-mixing mechanism: the reported checkpoints are not accompanied here by matched block-allocation measurements or mixing interventions.

\paragraph{Targeted adaptation and intrinsic retention.}
The protocol constructs an error set from TriviaQA/HotpotQA training examples that the frozen model fails under sampled generation, and trains on those examples. Post-training accuracy on this set measures learning of the targeted training examples; it is not held-out adaptation accuracy. Retention is evaluated on non-trained examples. The analysis includes examples with nonzero pre-adaptation answer scores and held-out validation examples; only examples with high initial scores can contribute to the severe negative-shift count below. With normalized pre/post scores $S_{\pre}(x),S_{\mathrm{post}}(x)$, that count uses
\begin{equation}
\mathbb I_{\mathrm{neg}}(x)=\mathbb I[S_{\mathrm{post}}(x)-S_{\pre}(x)<-0.8].
\end{equation}
The result-processing script uses the strict inequality. For scores that are exact fractions of ten samples, this corresponds to a loss of more than eight successful answers. It measures severe forgetting rather than every decline in answer probability.

\paragraph{Evaluation and training settings.}
The QA evaluation protocol is 5-shot with ten sampled answers per prompt, temperature 0.5, top-$p$ 0.9, and a 20-token generation limit, using answer-alias matching. The broad evaluation includes 0-shot MMLU~\citep{hendrycks2020measuring} and a nine-task BIG-bench multiple-choice subset~\citep{srivastava2023beyond}. The latter comprises analogical similarity, analytic entailment, causal judgment, cause and effect, common morpheme, conceptual combinations, logical deduction, logical sequence, and odd one out. It is labeled \emph{Reasoning} in Table~\ref{tab:combined_model_performance}; it is not the full BIG-Bench Hard benchmark. The recorded training protocol uses SFTTrainer, ten epochs, fused AdamW, cosine scheduling with warmup ratio 0.05, gradient clipping at 1.0, per-device batch size 64, maximum sequence length 1024, and bfloat16 precision. Baseline adapters cover attention and feed-forward linear projections; reported sweeps use learning rates from $10^{-6}$ to $10^{-2}$, nominal scaling 32, no dropout, and global seed 42.

\paragraph{Interpretation.}
Table~\ref{tab:combined_model_performance} reports adaptation scores and forgetting counts across coarse capacity tiers: small ($\le10$M trainable parameters), medium ($10$M--$100$M), and large ($>100$M). These are not exact capacity matches, and scores are descriptive outcomes of configuration sweeps rather than matched spectral interventions. The comparisons show no uniform ordering across models, datasets, and capacity tiers. They also omit UILinLoRA, preventing attribution of an outcome to rotations rather than scalers. Their role is to distinguish functional evaluation from structural guarantees, not to infer retained knowledge from a weight-space constraint.

% Inactive modular snapshot; edit the root neurips_2026.tex instead.
\begin{table*}[t]
\centering
\small
\setlength{\tabcolsep}{5pt}
\begin{adjustbox}{max width=\linewidth}
\begin{tabular}{l | c ccc | c ccc}
\toprule
& \multicolumn{4}{c|}{\textbf{TriviaQA}} & \multicolumn{4}{c}{\textbf{HotpotQA}} \\
& \multicolumn{4}{c|}{\textit{(Train $N=32{,}799$ | Known $N=43{,}724$)}} & \multicolumn{4}{c}{\textit{(Train $N=10{,}938$ | Known $N=4{,}723$)}} \\
\cmidrule(lr){2-5} \cmidrule(lr){6-9}
\multirow{2}{*}{\textbf{Method}} & \textbf{Adaptation} & \multicolumn{3}{c|}{\textbf{Retention}} & \textbf{Adaptation} & \multicolumn{3}{c}{\textbf{Retention}} \\
\cmidrule(lr){2-2} \cmidrule(lr){3-5} \cmidrule(lr){6-6} \cmidrule(lr){7-9}
& \textbf{Acc} $\uparrow$ & \textbf{Forget} $\downarrow$ & \textbf{MMLU} $\uparrow$ & \textbf{Reasoning} $\uparrow$ & \textbf{Acc} $\uparrow$ & \textbf{Forget} $\downarrow$ & \textbf{MMLU} $\uparrow$ & \textbf{Reasoning} $\uparrow$ \\
\midrule
\multicolumn{1}{c}{} & \multicolumn{8}{c}{\textbf{Gemma-3-12B-it}} \\
\midrule
\addlinespace[0.3ex]
\rowcolor{gray!10} \textit{Base Model} & \textit{--} & \textit{0} & \textit{0.715} & \textit{0.638} & \textit{--} & \textit{0} & \textit{0.715} & \textit{0.638} \\
\midrule
\multicolumn{9}{l}{\textit{Small}} \\
VeRA         & 0.976 & 8,463 & 0.576 & 0.425 & 0.993 & 1,632 & 0.661 & 0.607 \\
UIOrthoLoRA & 0.982 & 8,522 & 0.696 & 0.582 & 0.985 & 1,547 & 0.678 & 0.629 \\
LoRA         & 0.996 & 8,920 & 0.549 & 0.453 & 0.990 & 1,589 & 0.690 & 0.618 \\
DoRA         & 0.997 & 8,825 & 0.571 & 0.458 & 0.989 & 1,570 & 0.604 & 0.487 \\
\addlinespace
\multicolumn{9}{l}{\textit{Medium}} \\
UIOrthoLoRA & 0.982 & 8,430 & 0.704 & 0.602 & 0.988 & 1,492 & 0.691 & 0.597 \\
LoRA         & 0.997 & 7,825 & 0.601 & 0.466 & 0.994 & 1,488 & 0.699 & 0.615 \\
DoRA         & 0.997 & 8,114 & 0.601 & 0.468 & 0.992 & 1,510 & 0.608 & 0.475 \\
\addlinespace
\multicolumn{9}{l}{\textit{Large}} \\
UIOrthoLoRA & 0.975 & 8,194 & 0.702 & 0.604 & 0.988 & 1,523 & 0.690 & 0.617 \\
LoRA         & 0.997 & 7,638 & 0.602 & 0.486 & 0.995 & 1,496 & 0.705 & 0.633 \\
DoRA         & 0.978 & 7,604 & 0.604 & 0.460 & 0.993 & 1,515 & 0.617 & 0.488 \\
\midrule
\multicolumn{1}{c}{} & \multicolumn{8}{c}{\textbf{Llama-3.2-3B-it}} \\
\midrule
\addlinespace[0.3ex]
\rowcolor{gray!10} \textit{Base Model} & \textit{--} & \textit{0} & \textit{0.622} & \textit{0.458} & \textit{--} & \textit{0} & \textit{0.622} & \textit{0.458} \\
\midrule
\multicolumn{9}{l}{\textit{Small}} \\
VeRA         & 0.925 & 13,805 & 0.420 & 0.415 & 0.904 & 1,350 & 0.250 & 0.353 \\
UIOrthoLoRA & 0.935 & 10,088 & 0.549 & 0.428 & 0.958 & 931 & 0.530 & 0.419 \\
LoRA         & 0.938    & 13,597    & 0.342    & 0.397    & 0.942 & 1,066 & 0.496 & 0.404 \\
DoRA         & 0.952 & 13,910 & 0.322 & 0.388 & 0.954 & 1,064 & 0.429 & 0.417 \\
\addlinespace
\multicolumn{9}{l}{\textit{Medium}} \\
UIOrthoLoRA & 0.945 & 10,234 & 0.544 & 0.440 & 0.960 & 906 & 0.522 & 0.455 \\
LoRA         & 0.961 & 8,967 & 0.601 & 0.466 & 0.993 & 971 & 0.548 & 0.459 \\
DoRA         & 0.963 & 8,908 & 0.601 & 0.468 & 0.994 & 982 & 0.527 & 0.461 \\
\addlinespace
\multicolumn{9}{l}{\textit{Large}} \\
UIOrthoLoRA & 0.943 & 10,108 & 0.559 & 0.416 & 0.960 & 870 & 0.549 & 0.449 \\
LoRA         & 0.966 & 8,703 & 0.601 & 0.456 & 0.994 & 892 & 0.581 & 0.462 \\
DoRA         & 0.968 & 8,556 & 0.603 & 0.467 & 0.996 & 925 & 0.570 & 0.455 \\
\bottomrule
\end{tabular}
\end{adjustbox}
\caption{Exploratory adaptation and retention across model and capacity settings. Adaptation is measured on targeted training examples; Forget counts severe negative shifts on non-trained examples. Reasoning denotes the nine-task BIG-bench subset specified in Appendix~\ref{app:retention}, not full BIG-Bench Hard. Values and partition sizes are retained from the reported study; the coarse capacity groups are not exact parameter matches. No mixing-regularized or UILinLoRA counterpart is included.}
\label{tab:combined_model_performance}
\end{table*}


\end{document}
